\documentclass[11pt]{article}
\usepackage{amsmath, amssymb}
\usepackage{geometry}
\geometry{margin=0.7in}
\setlength{\parindent}{0pt}
\setlength{\parskip}{0.35em}

\begin{document}
\begin{center}
{\LARGE Statistics and Probability Cheatsheet}\\[0.4em]
{\large Fundamental Review}
\end{center}

\section*{Measures of Center}
\[
\text{mean}=\frac{\text{sum of values}}{\text{number of values}}
\]
\[
\text{median}=\text{middle value of ordered data}
\]
\[
\text{mode}=\text{most frequent value}
\]

\section*{Measures of Spread}
\[
\text{range}=\text{max}-\text{min}
\]
\[
\mathrm{IQR}=Q_3-Q_1
\]

\section*{Five-Number Summary}
\[
\min,\quad Q_1,\quad \text{median},\quad Q_3,\quad \max
\]
This is what you need for a box-and-whisker plot.

\section*{Box Plot Notes}
\begin{itemize}
  \item box runs from $Q_1$ to $Q_3$
  \item line inside box is median
  \item whiskers extend to min and max
\end{itemize}

\section*{Counting Principle}
If one choice has $a$ options and another has $b$ options, total outcomes:
\[
ab
\]
Extend by multiplying more stages.

\section*{Basic Probability}
\[
P(\text{event})=\frac{\text{favorable outcomes}}{\text{total outcomes}}
\]

\section*{Complement}
\[
P(\text{not }A)=1-P(A)
\]

\section*{Independent Events}
\[
P(A\text{ and }B)=P(A)P(B)
\]

\section*{Simple Odds Language}
\begin{itemize}
  \item probability compares favorable to total
  \item odds in favor compare favorable to unfavorable
\end{itemize}

\end{document}
